\subsection{Лекция 17 (28.02.26) Задача об обтекании цилиндра. PISO + МКО}
В тесте \cvar{cylinder-fvm} из файла \ename{cylinder_fvm_test.cpp}
реализовано решение задачи о нестационарном обтекании цилиндра методом PISO 
с испльзованием конечнообъемной аппроксимации на неструктурированной сетке (см. \secref{sec:fvm_piso_steps}).
\paragraph{Постановка задачи}
представлена на \figref{fig:cylinder-problem}.
\begin{figure}[h!]
\centering
\includegraphics[width=0.8\linewidth]{cylinder_problem.pdf}
\caption{К постановке задачи об обтекании цилиндра}
\label{fig:cylinder-problem}
\end{figure}

\paragraph{Баланс сил. Коэффициенты сил. Сопротивление.}
Проинтегрируем уравнение движения 
\cref{eq:ns2d_u}
по области расчёта $D$:
$$
\arint{\dfr{u^2}{x}}{D}{\vec x} +
\arint{\dfr{uv}{y}}{D}{\vec x} =
-\arint{\dfr{p}{x}}{D}{\vec x} + \frac{1}{\Ren} \arint{\nabla^2 u}{D}{\vec x}.
$$
Интегрирование по частям даёт:
\begin{align*}
&\arint{\dfr{f}{x}}{D}{\vec x} = \arint{f}{\Gamma_{out}}{s} - \arint{f}{\Gamma_{in}}{s} + \arint{f n_x}{\gamma}{s} \\[10pt]
&\arint{\dfr{f}{y}}{D}{\vec x} = \arint{f}{\Gamma_{top}}{s} - \arint{f}{\Gamma_{bot}}{s}+ \arint{f n_y}{\gamma}{s}, \\[10pt]
&\arint{\nabla^2 f}{D}{\vec x} = \arint{\dfr{f}{n}}{\Gamma}{s} + \arint{\dfr{f}{n}}{\gamma}{s}, \quad \Gamma = \Gamma_{in} \cap \Gamma_{out} \cap \Gamma_{bot} \cap \Gamma_{top}
\end{align*}
Учтём, что на обтекаемом теле скорости равны нулю, а верхняя и нижняя границы непротекаемы.
Тогда
$$
\arint{\left(u^2 + p\right)}{\Gamma_{in}}{s}
-\arint{\left(u^2 + p\right)}{\Gamma_{out}}{s} =
\arint{p \, n_x}{\gamma}{s} - \frac1\Ren\arint{\dfr{u}{n}}{\gamma}{s}
$$
Полученное выражение есть баланс сил в $x$ направлении.
Слева стоит сила, обусловленная перепадом динамического давления
(если считать профили скорости на входе и на выходе примерно одинаковыми, то останется только перепад статического давления).
А справа - силы сопротивления потоку вследствии наличия препятствия. И эти силы уравнавешивают друг друга.
Первое слагаемое в правой части -- есть сопротивление формы, второе -- сопротивление трения из-за эффектов вязкости.

Коэффициенты этих сил имеют следующее выражение:
\begin{equation}
\label{eq:ns2d_cx}
\begin{array}{lll}
	C^p_x =& 2 \displaystyle\arint{p\,n_x}{\gamma}{s} & \text{-- коэффициент сопротивления формы} \\[20pt]
	C^f_x =& -\dfrac2\Ren \displaystyle\arint{\dfr{u}{n}}{\gamma}{s} & \text{-- коэффициент сопротивления трения} \\[20pt]
	C_x =& C^p_x + C^f_x & \text{-- коэффициент сопротивления}
\end{array}
\end{equation}
Чтобы из этих безразмерных коэффициентов получить реальные силы, измеряемые в Ньютонах, нужно умножить их на $\tfrac12 \rho U^2 L^2$.
\paragraph{Подъёмная сила}
Аналогично проинтегрируем уравнение движение в направлении $y$
\cref{eq:ns2d_v}. С учётом граничных условий получим выражение для баланса сил в поперечном направлении
$$
\arint{p}{\Gamma_{bot}}{s}
-\arint{p}{\Gamma_{top}}{s} =
\arint{p \, n_y}{\gamma}{s} - \frac1\Ren\arint{\dfr{v}{n}}{\gamma}{s}
$$
и соответствующие коэффициенты
\begin{equation}
\label{eq:ns2d_cy}
\begin{array}{lll}
	C^p_y =& 2 \displaystyle\arint{p n_y}{\gamma}{s} & \text{} \\[20pt]
	C^f_y =& -\dfrac2\Ren \displaystyle\arint{\dfr{v}{n}}{\gamma}{s} & \text{} \\[20pt]
	C_y =& C^p_y + C^f_y & \text{-- коэффициент подъёмной силы}
\end{array}
\end{equation}

Для работы необходимо увеличить конечное время расчёта \cvar{end_time}
до значения, на котором устанавливается автоколебательный режим течения.
Расчёты и замеры времени исполнения проводить в релизной сборке.
\begin{enumerate}
\item
Нарисовать поля течения.
Провести расчёт без аппроксимации Rhie-Chow и визуализировать шахматные осцилляции.
\item 
Посчитать коэффициенты сопротивления и подъемной силы цилиндра.
Нарисовать графики $C_x(t)$, $C_y(t)$.
\item
Реализовать метод PIMPLE (см. \secref{sec:fvm_pimple_steps}). Подобрать оптимальные значения шага по времени и количества
внутренних PISO итераций, а также параметров релаксаций $\alpha_u$, $\alpha_p$ для максимально быстрого расчёта.
Сравнить полученные графики $C_x(t)$, $C_y(t)$ для различных настроек.
\item
Замерить время исполнения отдельных шагов расчёта. Оптимизировать код для ускорения счёта.
\end{enumerate}

\paragraph{Рекомендации к программированию}
\begin{itemize}
\item п.1: Аппроксимация граничных значение пробной скорости $\unij^*$
осущестляется по методике \secref{sec:rhie_chow}
в методе \cvar{compute_un_face_rhie_chow}.
Чтобы отключить эту интерполяцию, при расчёте $\unij^*$
следует воспользоваться реализацией
\cvar{compute_un_face}
испльзующий наивный подход \cref{eq:ns2d_fvm_utilde_n}.

\item п.2: Для вычисления коэффициентов сил по формулам \cref{eq:ns2d_cx,eq:ns2d_cy}
           требуется определить лежащие на цилиндре точки коллокации и вычилсить в них производную скорости, давление и нормаль.
\begin{cppcode}
for (size_t icolloc: boundary_info_.cyl){   // цикл по точка коллокации на цилиндре
    size_t iface = collocations_.face_index(icolloc);  // индекс граничной грани
    double pbnd = p_[icolloc]; // давление на грани
    Vector n = grid_.face_normal(iface); // нормаль к грани.
    double dudn = df_dn_computer_.compute(iface, u_); // dudn
    double dvdn = df_dn_computer_.compute(iface, v_); // dvdn
    if (grid_.tab_face_cell(iface)[0] == INVALID_INDEX){
        // сеточная нормаль смотрит внутрь. Развернём все направленные данные
        n *= -1;
        dudn *= -1;
        dvdn *= -1;
    }
    double face_area = grid_.face_area(iface);
    // ... вычисление коэффициентов
}
\end{cppcode}
\item п.3: Для реализации PIMPLE нужно вставить цикл по итерациям внутрь цикла по времени.
\begin{cppcode}
    for (time = time_step; time < end_time+1e-6; time += time_step){
        std::cout << "time= " << time << std::endl;
        worker.init_time_step();
        for (size_t it=0; it < max_it; ++it){    // <- SIMPLE
            double nrm = worker.init_step();
            worker.step(max_piso_iters, piso_eps);
        }
        worker.save_current_fields(time);
    }
\end{cppcode}
Метод \cvar{init_step()} возвращает невязку текущих значений решения. Её можно
использовать для контроля сходимости.

\item п.4: Для улучшения производительности в первую очередь следует сменить решатель
           для уравнений моментов:
\begin{cppcode}
AmgcMatrixSolver::solve_slae(mat_u_, rhs_u_, u_star);
AmgcMatrixSolver::solve_slae(mat_v_, rhs_v_, v_star);
\end{cppcode}
           Вместо многосеточного решателя предлагается испльзовать простые
           итерационные алгоритмы (\cref{sec:sae_simple}). Например SOR.
           При этом доводить их до сходимости не обязательно. Так как за окончательную сходимость отвечает внешний SIMPLE цикл.

           Также стоит рассмотреть возможность сборки диффузной части матрицы A (из метода
           \cvar{assemble_u_matrix()})  один раз на шаге инициализации.
           В самом методе вместо вычисления диффузии просто прибавить предрасчитанную матрицу.
           Если у предрасчитанной диффузной и итоговой матрицы одинаковый шаблон, то можно просто добавить val массив csr представления.
\end{itemize}
